\begin{abstract}
The Open Container Initiative (OCI) decouples image format, runtime, and tooling so that the same image can be executed by different runtimes and orchestrated by control planes. This modularity, together with the open-source nature of the reference runtime \texttt{runc}, makes the runtime layer itself a viable place to enforce security defaults, rather than relying on operators to supply hand-written seccomp, AppArmor, and capability profiles for every workload.

This thesis presents a security-first fork of \texttt{runc} that (i)~ships with a minimal default capability set instead of the historical Docker fourteen, (ii)~adds a single \texttt{-{}-security-scan} flag that, during a regular \texttt{runc run}, traces the workload via eBPF and emits cgroup-wide capability traces, a syscall allow-list, and an AppArmor complain-mode profile, and (iii)~narrows the on-disk \texttt{config.json} to exactly the privileges that the workload actually exercised. All hooks are shipped inside the runtime binary itself via hidden \texttt{scan-*} sub-commands, so the OCI bundle never has to embed executable artefacts.

The runtime is situated against manual policy authoring, cluster-side recorders such as the Kubernetes Security Profiles Operator and Inspektor Gadget, and---analytically, from the literature---sandboxed runtimes (gVisor, Kata Containers). The central claim is that a developer with one CLI flag and no cluster prerequisites can reach confinement comparable to what a skilled operator would produce by hand or what a staffed CI pipeline would record, while teams that cannot afford orchestration infrastructure remain covered.

Evaluation combines two strands: an empirical security laboratory (\texttt{runc-hardened-test}) that compares stock \texttt{runc} with scan-then-enforce on an intentionally vulnerable workload, and a qualitative analysis grounded in the literature and the competitive landscape summarised in Chapter~\ref{ch:litreview}. Performance benchmarks against sandbox runtimes are left to future work.
\end{abstract}
